\section{拓扑哈密顿量及其流的唯一性}
\label{sec:uniqueness-topological-hamiltonian-flow}

本节研究拓扑哈密顿量及其流的基本唯一性定理。

流的唯一性可以重新表述为展开映射
\[
 \overline\pi_2:\overline{\mathcal Q}
 \longrightarrow L_0^{(1,\infty)}([0,1]\times M,\mathbb R)
\]
的单射性。

\begin{theorem}\label{thm:uniqueness-topological-hamiltonian-flow}
设
$\lambda_i=\phi_{H_i}$、$\lambda_i'=\phi_{H_i'}\in
\mathcal P^{\mathrm{ham}}(\Symp(M,\omega),\operatorname{id})$
是两个光滑哈密顿路径序列，并且
\begin{enumerate}[label=(\arabic*)]
 \item $H_i$ 与 $H_i'$ 都在 $L^{(1,\infty)}$-拓扑中收敛，且
 $\|\overline H_i\#H_i'\|\to0$；
 \item $\phi_{H_i}$ 与 $\phi_{H_i'}$ 都在 $[0,1]\times M$ 上一致收敛。
\end{enumerate}
则
\[
 \lim_i\phi_{H_i}=\lim_i\phi_{H_i'}.
\]
\end{theorem}

\begin{proof}
先注意，由于 $\phi_{H_i}$ 与 $\phi_{H_i'}$ 都一致收敛，结论等价于
\[
 \lim_i(\phi_{H_i}^t)^{-1}\phi_{H_i'}^t\longrightarrow
 \operatorname{id}
\]
当 $i\to\infty$ 时关于 $t\in[0,1]$ 一致收敛。因此，考虑
$\overline H_i\#H_i'$ 后，证明归结为如下特殊情形。假设序列
$H_i$ 满足
\begin{enumerate}[label=(\arabic*)]
 \item $\|H_i\|\to0$；
 \item $\phi_{H_i}$ 在 $[0,1]\times M$ 上一致收敛。
\end{enumerate}
则 $\phi_{H_i}$ 在 $C^0$-拓扑中收敛于常值恒等路径。

在证明余下部分，我们证明这个命题。记
$\lambda=\lim_{i\to\infty}\phi_{H_i}\in
\mathcal P(\operatorname{Homeo}(M),\operatorname{id})$，并考虑它的值
$\lambda(t)$。我们要证明对所有 $t$ 都有
$\lambda(t)\equiv\operatorname{id}$。反设不然，并考虑
\[
 t_0=\sup\{t\in[0,1]\mid
 \lambda(s)=\operatorname{id}\text{ 对所有 }s\leq t\}.
\]
我们要证明 $t_0=1$。反设 $t_0<1$。由于对某个
$t_1>0$ 有 $\lambda(t_1)\ne\operatorname{id}$，而
$\lambda(0)=\operatorname{id}$，故知 $t_0>0$。于是可以选取充分小的
$\delta>0$，使得 $0<t_0+\delta<1$ 且
$\psi=\lambda(t_0+\delta)\ne\operatorname{id}$。注意，$\psi$ 必须连续，因为它是连续映射
$\phi_{H_i}^{t_0+\delta}$ 的一致极限。于是可以找到一个小闭球 $B$，使得
\[
 B\cap\psi(B)=\varnothing.
\]
由于 $B$，从而 $\psi(B)$ 是紧的，并且
$\phi_{H_i}^{t_0+\delta}\to\psi$ 一致收敛，对所有充分大的 $i$ 都有
\[
 B\cap\bigl(\phi_{H_i}^{t_0+\delta}\bigr)(B)=\varnothing.
\]
由霍费尔位移能 $e$ 的定义，有
$e(B)\leq\|\phi_{H_i}^{t_0+\delta}\|$。由定理 5.4.4 的能量--容量不等式，知道
$e(B)>0$，从而对所有充分大的 $i$，
\[
 0<e(B)\leq\|H_i^{t_0+\delta}\|.
\]
这里 $H_i^u$ 是生成 $\phi_{H_i}$ 的重参数化子路径
\[
 t\in[0,1]\longmapsto\phi_{H_i}^{ut}
\]
的哈密顿量；由显式公式，已经知道
$H_i^u(t,x)=uH_i(ut,x)$。因此，对于 $u=t_0+\delta$，已经导出
\[
\begin{aligned}
 0<e(B)
 &\leq\int_0^1\operatorname{osc}(H_{i,t}^u)\,dt
 =\int_0^1u(\max H_{i,ut}-\min H_{i,ut})\,dt\\
 &=\int_0^u(\max H_{i,t}-\min H_{i,t})\,dt
 \leq\|H_i\|\longrightarrow0,
\end{aligned}
\]
这给出矛盾。因此 $t_0=1$，证明完毕。
\end{proof}

接着考察 \eqref{eq:projection-topological-hamiltonian-paths} 中映射
$\overline\pi_1:\overline{\mathcal Q}\to
\mathcal P(\operatorname{Homeo}(M),\operatorname{id})$ 的单射性；我们把它表述为“拓扑哈密顿量的唯一性”。

\begin{question}\label{ques:uniqueness-limits-topological-hamiltonians}
假设 $H_i,H_i':[0,1]\times M\to\mathbb R$ 是两个满足归一化条件的哈密顿量序列。假设哈密顿路径
$\phi_{H_i},\phi_{H_i'}$ 一致收敛并且满足
\[
 d(\phi_{H_i},\phi_{H_i'})\longrightarrow0.
\]
那么，作为 $L^{(1,\infty)}$-函数，是否有
$\lim_{i\to\infty}H_i=\lim_{i\to\infty}H_i'$？
\end{question}

这个问题是在 (OhM07) 中提出的，下面的定理对它给出了肯定回答。

\begin{theorem}[维泰博，布霍夫斯基--赛法迪尼]
\label{thm:viterbo-buhovsky-seyfaddini-uniqueness}
\stfhIndex{I066}{bu^huo^fu^si^ji^sai^fa^di^ni}{布霍夫斯基--赛法迪尼（Buhovsky–Seyfaddini）}
假设 $H_i:[0,1]\times M\to\mathbb R$ 是一个哈密顿量序列。则在下列条件下，
$H_i\to h(t)$ 在 $L^{(1,\infty)}$-拓扑中收敛，其中
$h:[0,1]\to[0,1]$ 是一个 $L^1$-函数：
\begin{enumerate}[label=(\arabic*)]
 \item $H_i\to h(t)$ 在 $L^{(1,\infty)}$-拓扑中收敛，其中
 $h:[0,1]\to[0,1]$ 是一个 $L^1$-函数；
 \item $\phi_{H_i}$ 一致收敛于常值路径 $\operatorname{id}$。
\end{enumerate}
特别地，如果 $H_i$ 均作了均值归一化，则 $h\equiv0$。
\end{theorem}

这个定理最初由维泰博 (Vi06) 在闭情形的 $C^0$ 背景（或
$L^\infty$ 背景）下证明，随后在 (Oh07) 中推广到了开流形上的紧支情形。这个最初的证明依赖于某种
$C^0$-拉格朗日交理论，它本身也十分有趣。

我们给出布霍夫斯基与赛法迪尼 (BS10) 的一个优美证明，它建立在对经典勒贝格微分定理的巧妙应用上。文献中的标准勒贝格微分定理是针对\emph{实值}函数陈述的。他们证明中的一个关键要素，是针对定义 6.2.5 意义下属于
$L^1([0,1],C^0(M))$ 的函数给出一个微分定理。我们陈述这种一般性下的勒贝格微分定理版本。

先回忆经典分析中的几个基本定义。

\begin{definition}\label{def:lebesgue-set}
设 $f\in L^1_{\mathrm{loc}}(\mathbb R^n)$。定义 $f$ 的\emph{勒贝格集} $L_f$ 为
\stfhIndex{I194}{lei^bei^ge^ji}{勒贝格集（Lebesgue set）}
\[
 L_f=\left\{x\in\mathbb R^n\mathrel{\bigg|}
 \lim_{r\to0}\frac1{m(B_r(x))}
 \int_{B_r(x)}|f(y)-f(x)|\,dy=0\right\},
\]
其中 $B_r(x)$ 是以 $x$ 为心、半径为 $r$ 的球，而
$m(B_r(x))$ 是 $B_r(x)$ 的勒贝格测度。
\end{definition}

下面是 $\mathbb R^n$ 上函数论中的经典微分定理，其证明可见任何标准实分析教材，例如 (Fol99)。

\begin{theorem}[勒贝格微分定理]\label{thm:lebesgue-differentiation}
\stfhIndex{I193}{lei^bei^ge^wei^fen^ding^li}{勒贝格微分定理（Lebesgue differentiation theorem）}
设 $f:\mathbb R^n\to\mathbb R$ 是一个几乎处处定义的实值函数。以
$(L_f)^c$ 表示 $L_f$ 在 $\mathbb R^n$ 中的补集。如果
$f\in L^1_{\mathrm{loc}}(\mathbb R^n;\mathbb R)$，则
$m((L_f)^c)=0$。
\end{theorem}

我们还需要更一般的微分定理。为此，先回忆 (Fol99) 中的如下概念，再陈述勒贝格微分定理的一个更一般形式。

\begin{definition}\label{def:shrinks-nicely}
称 $\mathbb R^n$ 的博雷尔子集族 $\{E_r\}_{r>0}$\emph{良好收缩到} $x\in\mathbb R^n$，如果
\begin{enumerate}[label=(\arabic*)]
 \item 对每个 $r$，都有 $E_r\subset B_r(x)$；
 \item 存在与 $r$ 无关的常数 $\alpha>0$，使得不等式
 $m(E_r)>\alpha m(B_r(x))$ 成立。
\end{enumerate}
这里集合 $E_r$ 本身不必包含 $x$。
\end{definition}

\begin{theorem}\label{thm:lebesgue-differentiation-nice-shrinking}
假设 $f\in L^1_{\mathrm{loc}}$。对 $f$ 的勒贝格集中的每个
$x$，特别是对几乎每个 $x$，都有
\begin{equation}
 \lim_{t\to0}\frac1{m(E_r)}
 \int_{E_r}|f(y)-f(x)|\,dy=0,
 \label{eq:lebesgue-differentiation-absolute}
\end{equation}
以及
\begin{equation}
 \lim_{t\to0}\frac1{m(E_r)}
 \int_{E_r}f(y)\,dy=f(x),
 \label{eq:lebesgue-differentiation-average}
\end{equation}
其中 $\{E_r\}_{r>0}$ 是任意一个良好收缩到 $x$ 的集合族。
\end{theorem}

对每个 $x\in M$ 以及
$H\in\mathcal H_0^{(1,\infty)}([0,1]\times M,\mathbb R)$，我们将把这个定理应用于
$L^1$-函数
\[
 t\in[0,1]\subset\mathbb R\longmapsto H(t,x)\in\mathbb R
\]
以及集合族
\[
 E_h(t):=[t,t+h],\qquad h>0,\quad t\in[0,1].
\]
注意，集合族 $\{E_h(t)\}_{h>0}$ 良好收缩到 $t$。下面陈述 (BS10) 中使用的微分定理版本。

\begin{proposition}[布霍夫斯基--赛法迪尼]
\label{prop:topological-hamiltonian-lebesgue-differentiation}
考虑
$H\in\mathcal H_0^{(1,\infty)}([0,1]\times M)
\subset L^{(1,\infty)}([0,1]\times M)$。则
\[
 \lim_{h\to0+}\frac1h\int_t^{t+h}
 \operatorname{osc}(H_s-H_t)\,ds=0
\]
对几乎每个 $t\in[0,1]$ 成立。
\end{proposition}

\begin{proof}
设 $H_i$ 是 $H$ 的一个逼近序列。则有收敛
\[
 \int_0^1\operatorname{osc}(H_i-H)\,dt\longrightarrow0,
\]
特别地，
\[
 \int_0^1\operatorname{osc}(H_i-H_j)\,dt\longrightarrow0
 \qquad\text{当 }i,j\to\infty.
\]
还有
\begin{equation}
\begin{split}
 \int_t^{t+h}\operatorname{osc}(H_s-H_t)\,ds
 &\leq\int_t^{t+h}\operatorname{osc}(H_{i,s}-H_{i,t})\,ds\\
 &\quad+\int_t^{t+h}\operatorname{osc}(H_s-H_{i,s})\,ds\\
 &\quad+\int_t^{t+h}\operatorname{osc}(H_{i,t}-H_t)\,ds.
\end{split}
 \label{eq:oscillation-three-term-bound}
\end{equation}
由于 $H_i$ 光滑，
\[
 \lim_{h\to0+}\frac1h\int_t^{t+h}
 \operatorname{osc}(H_{i,s}-H_{i,t})\,ds=0
\]
对所有 $t\in[0,1]$ 成立。另一方面，正如已经在定理 6.2.6 (2) 中证明的，函数
$\operatorname{osc}(H)$ 属于 $L^1([0,1],\mathbb R)$，并且
$\operatorname{osc}(H-H_i)\to0$ 于
$L^1([0,1],\mathbb R)$。记 $F_i=H-H_i$，并对每个 $i$ 把定理
6.3.5 应用于 $f_i:=\operatorname{osc}(F_i)$，从而证明
\[
 \lim_{h\to0+}\frac1h\int_t^{t+h}
 \bigl|\operatorname{osc}(F_{i,s})-
 \operatorname{osc}(F_{i,t})\bigr|\,ds=0
\]
对几乎每个 $t$ 成立。以 $L_{f_i}$ 表示 $f_i$ 的勒贝格集，并考虑集合
\[
 \bigcap_{i=1}^{\infty}L_{f_i}=:A\subset[0,1].
\]
这个集合在 $[0,1]$ 中具有满测度。又因为 $f_i\to0$ 于
$L^1$，所以它在 $[0,1]$ 上几乎处处逐点收敛于 $0$；设这发生在满测度子集
$B\subset[0,1]$ 上。

我们断言，在每个 $t\in A\cap B$ 处，
\begin{equation}
 \lim_{h\to0+}\int_t^{t+h}
 \operatorname{osc}(H_s-H_t)\,ds=0.
 \label{eq:oscillation-lebesgue-claim}
\end{equation}
在每个 $t\in A$ 处，有
\[
\begin{aligned}
 \lim_{h\to0+}\frac1h\int_t^{t+h}\operatorname{osc}(F_{i,s})\,ds
 &\leq\lim_{h\to0+}\frac1h\int_t^{t+h}
 \bigl|\operatorname{osc}(F_{i,s})-\operatorname{osc}(F_{i,t})\bigr|\,ds\\
 &\quad+\lim_{h\to0+}\frac1h\int_t^{t+h}
 \operatorname{osc}(F_{i,t})\,ds\\
 &=\lim_{h\to0+}\frac1h\int_t^{t+h}|f_i(s)-f_i(t)|\,ds+f_i(t)\\
 &=f_i(t).
\end{aligned}
\]
把它代入 \eqref{eq:oscillation-three-term-bound}，得到
\[
 \lim_{h\to0+}\int_t^{t+h}
 \operatorname{osc}(H_s-H_t)\,ds\leq f_i(t)
\]
对所有 $i$ 成立。由于 $t\in B$，有 $f_i(t)\to0$。令
$i\to\infty$，得到
\[
 \lim_{h\to0+}\int_t^{t+h}
 \operatorname{osc}(H_s-H_t)\,ds=0.
\]
因为 $A\cap B$ 具有满测度，证明完毕。
\end{proof}

现在可以证明定理 6.3.3 了。(BS10) 的一个关键思想在于研究\emph{零哈密顿量}集合
\stfhIndex{I217}{ling^ha^mi^dun^liang}{零哈密顿量（null Hamiltonian）}
\[
 \mathcal H_{\mathrm{null}}
 =\bigl\{H\in\mathcal H_0^{(1,\infty)}([0,1]\times M,\mathbb R)
 \mid\phi_H\equiv\operatorname{id}\bigr\}
 =(\overline\iota_{\mathrm{ham}})^{-1}(\operatorname{id})
\]
及其自治对应物
\[
 \mathcal H_{\mathrm{null}}^{\mathrm{aut}}
 =\{H\in\mathcal H_{\mathrm{null}}\mid H\text{ 与时间无关}\}
 \subset C^0(M).
\]
唯一性可以重新表述为 $\mathcal H_{\mathrm{null}}=\{0\}$。本节余下部分将证明这个命题。

证明由一系列引理与练习组成。

\begin{lemma}\label{lem:smooth-null-hamiltonian-zero}
如果
$H\in\mathcal H_{\mathrm{null}}\cap C^\infty([0,1]\times M)$，则
$H\equiv0$。
\end{lemma}

\begin{proof}
设 $H_i$ 是一个表示 $H$ 的光滑哈密顿量序列，即
$\|H_i-H\|\to0$ 且 $\phi_{H_i}\to\operatorname{id}$。由定理
6.3.1，有 $\phi_H\equiv\operatorname{id}$。对流
$\phi_H^t$ 微分，得到 $X_H\equiv0$，继而得到
$dH_t=0$，所以对所有 $t$，$H_t$ 都是常数。由归一化条件，这蕴含对所有
$t$ 都有 $H_t\equiv0$。
\end{proof}

\begin{exercise}\label{ex:null-hamiltonian-properties}
证明 $\mathcal H_{\mathrm{null}}$ 的下列性质：
\begin{enumerate}[label=(\arabic*)]
 \item $\mathcal H_{\mathrm{null}}$ 是 $L^{(1,\infty)}$ 的闭子集，但不一定是向量子空间；
 \item 对任意光滑递增函数 $\rho:[0,1]\to\mathbb R$，
 $\mathcal H_{\mathrm{null}}$ 在运算 $H\mapsto H^\rho$ 下封闭，其中
 $H^\rho(t,x)=\rho'(t)H(\rho(t),x)$；
 \item $\mathcal H_{\mathrm{null}}$ 对取和与取负运算封闭。
\end{enumerate}
\end{exercise}

\begin{lemma}\label{lem:autonomous-null-hamiltonians-vector-space}
$\mathcal H_{\mathrm{null}}^{\mathrm{aut}}$ 是 $\mathbb R$ 上的向量空间，并且
$\Symp(M,\omega)$ 在 $\mathcal H_{\mathrm{null}}^{\mathrm{aut}}$ 上的作用
\[
 h\longmapsto\psi^*h=h\circ\psi
\]
保持 $\mathcal H_{\mathrm{null}}^{\mathrm{aut}}$。
\end{lemma}

\begin{proof}
先证明第二个命题。设 $H_i$ 是表示
$h\in\mathcal H_{\mathrm{null}}^{\mathrm{aut}}$ 的一个序列。于是哈密顿量
$H_i(t,\psi(x))$ 表示 $h\circ\psi$，因为
$H_i(t,\psi(x))$ 生成流
$\psi\circ\phi_{H_i}^t\circ\psi^{-1}$。由于
$\phi_{H_i}\to\operatorname{id}$，得到
$\psi\circ\phi_{H_i}^t\circ\psi^{-1}\to\operatorname{id}$。所以
$h\circ\psi\in\mathcal H_{\mathrm{null}}^{\mathrm{aut}}$。

现在证明第一个命题。显然
$0\in\mathcal H_{\mathrm{null}}^{\mathrm{aut}}$。设
$h,k\in\mathcal H_{\mathrm{null}}^{\mathrm{aut}}$，并分别以
$H_i,K_i$ 表示它们的相伴逼近序列。于是哈密顿量
$H_i\#K_i$ 表示流 $\phi_{H_i}\circ\phi_{K_i}$。由于
$\phi_{H_i},\phi_{K_i}\to\operatorname{id}$，故
$\phi_{H_i}\circ\phi_{K_i}\to\operatorname{id}$。另一方面，依定义有
\[
 (H_i\#K_i)(t,x)
 =H_i(t,x)+K_i\bigl(t,(\phi_{H_i}^t)^{-1}(x)\bigr).
\]
由于 $H_i\to h$、$K_i\to k$ 于 $L^{(1,\infty)}$，便得到
$H_i\#K_i\to h+k$。

\begin{exercise}\label{ex:composition-hamiltonian-convergence}
证明函数
\[
 (t,x)\longmapsto K_i\bigl(t,(\phi_{H_i}^t)^{-1}(x)\bigr)
\]
在 $L^{(1,\infty)}$ 中收敛于 $k$。
\end{exercise}

类似地，回忆
\[
 \overline H_i(t,x)=-H_i(t,\phi_{H_i}(x))
\]
表示逆流 $\phi_{H_i}^{-1}$，便证明了
$-h\in\mathcal H_{\mathrm{null}}^{\mathrm{aut}}$。由于
$\phi_{H_i}\to\operatorname{id}$，有
$\phi_{H_i}^{-1}\to\operatorname{id}$，又有 $H_i\to h$ 以及
$-H_i(t,\phi_{H_i}(x))\to-h$ 于 $L^{(1,\infty)}$。证明完毕。
\end{proof}

\begin{lemma}\label{lem:null-hamiltonian-time-slice-autonomous}
设 $H\in\mathcal H_{\mathrm{null}}$。那么，对使得
$H_t\in C^0(M)$ 的几乎每个 $t\in[0,1]$，含时无关函数
\[
 h(x):=H_t(x)=H(t,x)
\]
属于 $\mathcal H_{\mathrm{null}}^{\mathrm{aut}}$。
\end{lemma}

\begin{proof}
由命题 6.3.8，对几乎每个 $t\in[0,1]$，都有
\[
 \lim_{\epsilon\to0+}\frac1\epsilon
 \int_t^{t+\epsilon}\|H_s-H_t\|_\infty\,ds=0.
\]
对充分大的 $N\in\mathbb N$，考虑函数
\[
 G_N(s,x)=\frac1N H\left(t+\frac sN,x\right),
\]
它对应于取 $\rho(s)=t+s/N$ 时的 $H^\rho$。因此
$G_N\in\mathcal H_{\mathrm{null}}$。由于练习 6.3.10 (3) 表明
$\mathcal H_{\mathrm{null}}$ 对取和运算封闭，还有
$H_N:=NG_N\in\mathcal H_{\mathrm{null}}$。但
$NG_N(s,x)=H(t+s/N,x)$。令 $\epsilon=1/N$，得到
\[
 \int_0^1\|H_N(s,\mathord\cdot)-h(\mathord\cdot)\|_\infty\,ds
 =\frac1\epsilon\int_t^{t+\epsilon}
 \|H(\tau,\mathord\cdot)-H(t,\mathord\cdot)\|_\infty\,d\tau
 \longrightarrow0,
\]
其中，对给定的固定 $t\in[0,1]$，函数 $h$ 定义为
$h(x)=H_t(x)$。所以当 $N\to\infty$ 时，
$H_N\to h$ 于 $L^{(1,\infty)}$。由于
$\mathcal H_{\mathrm{null}}$ 是 $L^{(1,\infty)}$ 的闭子集，故
$h\in\mathcal H_{\mathrm{null}}$；又因它与时间无关，继而得到
$h\in\mathcal H_{\mathrm{null}}^{\mathrm{aut}}$。
\end{proof}

下面是布霍夫斯基与赛法迪尼的另一个关键引理，它以一种新颖方式使用了经典分析中的工具。

\begin{lemma}\label{lem:autonomous-null-hamiltonian-zero}
如果 $h\in\mathcal H_{\mathrm{null}}^{\mathrm{aut}}$，则 $h=0$。
\end{lemma}

\begin{proof}
用反证法。假设 $h\ne0$。在这个假设下，我们将构造一个属于
$\mathcal H_{\mathrm{null}}^{\mathrm{aut}}\subset
\mathcal H_{\mathrm{null}}$ 的非零光滑函数，从而与引理 6.3.9 矛盾。

由于 $h$ 不是常数且 $M$ 连通，$h$ 不是局部常数。因此存在
$x_0\in M$，使得 $h$ 在 $x_0$ 的任何开邻域上都不是常数。选取
$x_0$ 的一个达布邻域 $U$。对任意满足
$h(y_0)\ne h(x_0)$ 的 $y_0\in U$，可以选取
$\phi\in\Ham_U(M,\omega)$，使得 $\phi(x_0)=y_0$。考虑函数
$k=h\circ\phi-h$。由 $y_0$ 与 $\phi$ 的选取，$k$ 是非常值函数，并且由引理
6.3.11，$k\in\mathcal H_{\mathrm{null}}^{\mathrm{aut}}$。此外，
$k=0$ 于 $\partial U$ 附近。由于 $k=0$ 于 $\partial U$ 附近，可以在
$U$ 外以 $0$ 把 $k$ 延拓到整个 $M$。定义 $\mathcal L$ 为 $M$ 上形如
\[
 \psi^*k=k\circ\psi,
 \qquad \operatorname{supp}k\subset U,\quad
 \psi\in\Ham_U(M,\omega)
\]
的所有函数之线性张成的 $C^0$-闭包。于是
$\mathcal L\subset\mathcal H_{\mathrm{null}}^{\mathrm{aut}}$。现在断言
$\mathcal L$ 包含一个非零光滑元素。

由于 $k=0$ 于 $U$ 外且 $U$ 是一个达布坐标卡，可以假设
$U\subset\mathbb R^{2n}$，而且只需在开集
$U\subset\mathbb R^{2n}$ 上证明这个命题。仍以 $\mathcal L$ 表示形如
$\psi^*k$ 的函数之线性张成的 $C^0$-闭包，其中
$\psi\in\Ham_U(\mathbb R^{2n})$。

对 $v\in\mathbb R^{2n}$，定义
$k_v(x)=k(x-v)$。接着证明，如果 $\|v\|$ 充分小，则
$k_v\in\mathcal L$。取 $\operatorname{supp}k$ 的一个邻域
$W$，使得 $\overline W\subset U$，并把 $v$ 取得如此之小，使得对所有
$0\leq t\leq1$，都有 $\overline W+tv\subset U$。选取截断函数
$\rho$，使得 $\operatorname{supp}\rho\subset U$，且
$\rho=1$ 于
$\bigcup_{t\in[0,1]}(\overline W+tv)$。接着考虑函数
\[
 g_v(x)=\omega(v,x)\rho.
\]
$g_v$ 的哈密顿流 $\phi_{g_v}^t$ 在
$\bigcup_{t\in[0,1]}(\overline W+tv)$ 上无非就是对所有
$x\in W$ 的流 $t\mapsto x+tv$，因而在
$\operatorname{supp}k$ 上也是如此。所以
\[
 \bigl((\phi_{g_v}^1)^{-1}\bigr)^*k
 =k\circ(\phi_{g_v}^1)^{-1}=k_v.
\]
这证明了 $k_v\in\mathcal L$。

现在选取一个光滑函数 $\chi$，其支集包含在
$\mathbb R^{2n}$ 中原点的一个充分小邻域内，并考虑由
\[
 k*\chi(x)=\int_{\mathbb R^{2n}}k(x-v)\chi(v)\,dv
\]
定义的卷积 $k*\chi$。我们断言 $k*\chi\in\mathcal L$：用黎曼和逼近该积分，可以把
$k*\chi$ 写成如下和的 $C^0$-极限：
\[
 \sum_{i=1}^m c_i k_{v_i}.
\]

\begin{exercise}\label{ex:convolution-c0-convergence}
证明这个 $C^0$-收敛。
\end{exercise}

于是，由 $\sum_{i=1}^m c_i k_{v_i}\in\mathcal L$ 及
$\mathcal L$ 的定义，得到 $k*\chi\in\mathcal L$。

最后，取狄拉克 $\delta_0$ 函数的磨光逼近 $\chi_\epsilon$，使得
\[
 k=\lim_{\epsilon\to0}k*\chi_\epsilon
 \qquad\text{于 }C^0.
\]
由于 $k\ne0$，当 $\epsilon$ 充分小时，$k*\chi_\epsilon$ 是非零光滑函数。由上述断言，每个
$k*\chi_\epsilon\in\mathcal L\subset
\mathcal H_{\mathrm{null}}^{\mathrm{aut}}$。这完成了引理
6.3.14 的证明。
\end{proof}

\begin{proof}[定理 6.3.3 证明的收束]
设 $H\in\mathcal H_{\mathrm{null}}$。那么，在
$H\in L^1([0,1],C^0(M))$ 的每个勒贝格点 $t$，有
$H_t\in\mathcal H_{\mathrm{null}}^{\mathrm{aut}}$，因而
$H_t\equiv0$。由于勒贝格点集合在 $[0,1]$ 中具有满测度，
$H=0$ 于 $L^1([0,1],C^0(M))$。所以
$H=0$ 于 $L^{(1,\infty)}([0,1]\times M,\mathbb R)$。证明完毕。
\end{proof}

对上述证明作一个小修改，还会导出拓扑哈密顿量的局部性。细节请读者参阅 (BS10)。这个局部性结果最初在 (Oh07) 中于
$C^0$ 背景下用另一种方法证明；该方法遵循维泰博证明的精神，使用关于开子集之余法丛的拉格朗日交定理。

\begin{theorem}[吴龙根 (Oh07)，布霍夫斯基--赛法迪尼 (BS10)]
\label{thm:locality-topological-hamiltonian}
假设 $H\in\mathcal H^{(1,\infty)}$，并且它的拓扑哈密顿流满足
$\phi_H^t|_U\equiv\operatorname{id}_U$。那么，在每个时刻 $t$，存在常数
$c(t)$，使得对所有 $x\in[0,1]\times U$ 都有
$H(t,x)=c(t)$。
\end{theorem}
